% !TEX root = thesis.tex
\chapter{Conclusion and Future Works}\label{sec:conclusion}

This work investigated the efficacy of two distinct learned representations—Large Wireless Models (LWM) and Autoencoders—for physical layer wireless communications. These representations were evaluated across three benchmark tasks, focusing on data efficiency, noise robustness, and computational complexity. Where possible, the underlying performance drivers were explained through geometric latent-space analysis.

Initially, the LWM showed promising capabilities. The transformer-based representation exhibited good data efficiency in the Line-of-Sight (LoS) and Non-Line-of-Sight (NLoS) classification task, despite exhibiting poor noise robustness. Conversely, it demonstrated lower data efficiency and noise robustness than the raw baseline in the beam selection task. In the final power allocation task, the model yielded mixed results, showing improved data efficiency and robustness only under specific, highly dense spatial scenarios.

Ultimately, the biggest challenge for the LWM representation is its immense computational cost. Across every task, utilizing the LWM incurred a severe latency and parameter overhead. In many cases, active fine-tuning was strictly necessary to make the CLS token representation competitive against standard baselines. For these reasons, the integration of such foundation models into real-time, physical layer applications must be weighed carefully against their hardware constraints.

In stark contrast, the compressed latent vector generated by the autoencoder demonstrated superior, and highly stable performance compared to both the LWM and the classical baselines. It achieved greater data efficiency and noise robustness while introducing only marginal computational overhead.

\begin{table}[htbp]
    \centering
    \caption{Qualitative performance comparison of learned representations against baselines.}
    \label{tab:overall_representation_comparison}
    \renewcommand{\arraystretch}{1.2} 
    \setlength{\tabcolsep}{4pt} 
    \small 
    \begin{tabular}{l l c c c}
        \hline
        \textbf{Task} & \textbf{Model} & \textbf{Data Eff.} & \textbf{Noise Robustness} & \textbf{Complexity} \\
        \hline
        LoS/NLoS & Autoencoder & Superior $\uparrow$ & Superior $\uparrow$ & Low Cost $\uparrow$ \\
        Classification & LWM & Good $\uparrow$ & Inferior $\downarrow$ & High Cost $\downarrow$ \\
        \hline
        Beam & Autoencoder & Superior $\uparrow$ & Superior $\uparrow$ & Low Cost $\uparrow$ \\
        Selection & LWM & Inferior $\downarrow$ & Inferior $\downarrow$ & High Cost $\downarrow$ \\
        \hline
        Power & Autoencoder & Superior $\uparrow$ & Superior $\uparrow$ & Low Cost $\uparrow$ \\
        Allocation & LWM & Mixed $\approx$ & Mixed$^*$ $\approx$ & High Cost $\downarrow$ \\
        \hline
    \end{tabular}
    
    \vspace{1ex}
    {\raggedright \footnotesize \textit{Note:} $(\uparrow)$ Better than baseline, $(\approx)$ Comparable, $(\downarrow)$ Worse than baseline. \\
    $^*$LWM showed a sum-rate spike at 16 users (high SNR), but underperformed in standard regimes.\par}
\end{table}

Based on these findings, we conclude that there is currently no practical advantage to expanding the dimensionality of the representation for these specific regression and classification tasks. In time-critical applications such as telecommunications, latency and computational cost are vital system constraints. The significant increase in memory footprint and inference time incurred by the LWM does not provide meaningful advantages that justify its deployment over simpler, compressed representations.

Furthermore, the low-dimensional nature of the autoencoder's output directly reduces the transmission overhead required for system feedback. In realistic telecommunication deployments, the User Equipment (UE) must transmit an estimation of the channel conditions back to the Base Station (BS). By feeding back the compressed latent vector $\mathbf{z}$, the BS is afforded the flexibility to either reconstruct the original channel state information or use $z$ directly to perform downstream tasks with minimal delay. The final comparison between different representations is shown in the \ref{tab:overall_representation_comparison}.

Despite these established advantages, there is still significant room for improvement in representational learning for wireless channels:

\begin{itemize}
    \item The current autoencoder architecture is heavily inspired by CSINet. It is imperative to research novel architectural paradigms that could further increase the compression ratio without sacrificing data efficiency.
    \item While the training methodology—based on reconstructing a clean channel from a noisy input—improved general noise resilience, it was not sufficient to make the model truly robust in ultra-low SNR scenarios. Advanced denoising objectives should be explored.
    \item Additionally, the pre-trained models utilized in this work were trained on a finite set of simulated environments. To improve generalization, it is necessary to drastically increase the complexity, scale, and diversity of the pre-training dataset.
\end{itemize}

The LWM served as a highly valuable initial experiment in applying foundation models to telecommunications. However, it is evident that directly porting techniques from Vision Transformers and Natural Language Processing is not sufficient to achieve optimal results in the physical layer. Future iterations of wireless foundation models could be improved through several key developments:

\begin{itemize}
    \item Designing a transformer architecture equipped with an inherent bottleneck to generate lower-dimensional representations, potentially merging the high-capacity learning of transformers with the computational efficiency of autoencoders.
    \item Developing a pre-training objective that explicitly accounts for physical noise distributions (AWGN, fading) in the input channel.
    \item Implementing a Mixture of Experts (MoE) architecture, allowing the foundation model to dynamically route and generate specialized representations for a diverse set of sub-tasks without inflating inference latency.
\end{itemize}

Finally, future research must expand the evaluation framework to include more complex, active downstream tasks—such as joint precoding and channel tracking—to comprehensively test the limits of these learned representations in fully dynamic 6G networks.